% !TeX root = ../main.tex
\section{Conclusion}
\label{sec:conclusion}

This paper presented a real-flight effective-wrench correction pipeline for a bird-scale flapping-wing vehicle. A DeLaurier-style aerodynamic model was used as a low-dimensional simulator prior, aligned with training logs through a force-channel affine wrapper, and then compared with six-component effective-wrench labels reconstructed from outdoor PX4 flight logs using measured mass, center-of-gravity, and inertia properties.

The main finding is that an affine-aligned flapping-wing aerodynamic prior leaves a structured real-flight force residual, and this residual can be reduced with a deployable phase- and condition-dependent correction. On held-out logs, the phase-structured correction reduced force-channel RMSE by $68.3\%$ for $f_{x,b}$, $15.0\%$ for $f_{y,b}$, and $69.8\%$ for $f_{z,b}$ relative to the affine-aligned prior. Effective moment channels were also predicted, with the strongest result in $m_{y,b}$, but the moment results remain more sensitive to inertia, reference-point consistency, and angular-acceleration smoothing.

Residual analysis showed that the prior mismatch is not arbitrary. The main force-channel residuals contain wingbeat-phase, frequency, and flight-condition structure, and the structured correction removes much of this repeatable component on held-out logs. This supports real-flight data as a practical way to build a log-conditioned correction around a low-dimensional flapping-wing simulator prior without replacing the prior with an unrestricted black-box model.

The corrected map is a prerequisite step toward simulator embedding, not a completed closed-loop simulation validation. Replay was used here only as a diagnostic because outdoor replay compounds estimator noise, smoothing choices, time alignment, wind disturbance, and controller-response mismatch. Future work will embed the correction into IsaacLab with explicit disturbance and uncertainty handling, evaluate controller-in-the-loop behavior, and collect broader flight data over angle-of-attack, flapping-frequency, sideslip, wind, and maneuver conditions.
